\documentclass[12pt,a4paper]{article}
%\usepackage[utf-8]{inputenc}
\usepackage[margin=1in]{geometry}
\usepackage{amsmath}
\usepackage{amssymb}
\usepackage{graphicx}
\usepackage{array}
\usepackage{booktabs}
\usepackage{xcolor}
\usepackage{fancyhdr}
\usepackage{hyperref}

\pagestyle{fancy}
\fancyhf{}
\rhead{GED Math Worksheet}
\lhead{Graphs and Averages}
\cfoot{\thepage}

\title{\textbf{GED Mathematics Worksheet}\\[0.3cm]\Large{Reading Graphs, Tables, and Computing Averages}}
\author{50 Practice Questions}
\date{May 2026}

\begin{document}

\maketitle

\tableofcontents

\newpage

\section{Instructions}
\begin{itemize}
    \item This worksheet contains 50 questions covering graphs, tables, and averages from GED Mathematics lessons.
    \item Questions are organized by topic to help you focus on specific skills.
    \item Show your work for all calculations.
    \item A solutions document with complete worked answers is available separately.
    \item Calculator use is permitted where appropriate.
\end{itemize}

\newpage

\section{Part A: Reading and Interpreting Tables (Questions 1-10)}

\begin{enumerate}

\item The table below shows quarterly sales (in thousands) for a retail store in 2025:
\begin{center}
\begin{tabular}{|c|c|c|c|c|}
\hline
Quarter & Q1 & Q2 & Q3 & Q4 \\
\hline
Sales & \$45 & \$62 & \$58 & \$71 \\
\hline
\end{tabular}
\end{center}
What was the total sales for the year?

\item Using the table from Question 1, by how much did sales increase from Q1 to Q4?

\item A table shows the number of hours worked by five employees in a week:
\begin{center}
\begin{tabular}{|c|c|c|c|c|c|}
\hline
Employee & Alice & Bob & Carol & David & Emma \\
\hline
Hours & 40 & 38 & 42 & 35 & 45 \\
\hline
\end{tabular}
\end{center}
Which employee worked the most hours?

\item Using the table from Question 3, how many total hours did all employees work?

\item A fitness center tracks daily membership sign-ups over one week:
\begin{center}
\begin{tabular}{|c|c|c|c|c|c|c|}
\hline
Day & Mon & Tue & Wed & Thu & Fri & Sat \\
\hline
Sign-ups & 8 & 12 & 15 & 11 & 18 & 22 \\
\hline
\end{tabular}
\end{center}
On which day did the fitness center have the lowest number of sign-ups?

\item Using the table from Question 5, what was the total number of sign-ups for the week?

\item A grocery store tracks weekly produce sales (in pounds):
\begin{center}
\begin{tabular}{|c|c|c|c|c|}
\hline
Produce & Apples & Broccoli & Carrots & Lettuce \\
\hline
Pounds Sold & 240 & 180 & 160 & 200 \\
\hline
\end{tabular}
\end{center}
How many more pounds of apples were sold than carrots?

\item Using the table from Question 7, if broccoli costs \$2.50 per pound wholesale, what is the total cost of broccoli sold?

\item A restaurant tracks customer satisfaction scores by meal type:
\begin{center}
\begin{tabular}{|c|c|c|c|c|}
\hline
Meal Type & Breakfast & Lunch & Dinner & Late Night \\
\hline
Avg. Score (out of 10) & 8.2 & 7.9 & 8.6 & 7.4 \\
\hline
\end{tabular}
\end{center}
Which meal type had the highest satisfaction score?

\item Using the table from Question 9, what is the difference between the highest and lowest satisfaction scores?

\end{enumerate}

\section{Part B: Box Plots and Data Interpretation (Questions 11-20)}

\begin{enumerate}
\setcounter{enumi}{10}

\item A box plot for test scores has the following values:
\begin{itemize}
    \item Minimum: 45
    \item Q1 (Lower Quartile): 62
    \item Median: 75
    \item Q3 (Upper Quartile): 85
    \item Maximum: 98
\end{itemize}
What is the interquartile range (IQR)?

\item Using the box plot from Question 11, how many of the data values fall between Q1 and Q3?

\item A dataset has the following five-number summary:
\begin{itemize}
    \item Min: 12
    \item Q1: 28
    \item Median: 42
    \item Q3: 56
    \item Max: 78
\end{itemize}
What is the range of the data?

\item Using the data from Question 13, what percentage of values are between the minimum and Q1?

\item A school tracks attendance percentages for a class throughout the year. The box plot shows:
\begin{itemize}
    \item Min: 78\%
    \item Q1: 88\%
    \item Median: 92\%
    \item Q3: 95\%
    \item Max: 99\%
\end{itemize}
Identify any outliers (values more than 1.5 times the IQR below Q1 or above Q3).

\item Using the attendance data from Question 15, how many quarters of the data lie above the median?

\item A box plot for monthly rainfall (in inches) shows:
\begin{itemize}
    \item Q1: 2.5
    \item Median: 4.0
    \item Q3: 5.5
\end{itemize}
What is the IQR for this data?

\item Using the rainfall data from Question 17, what is the lower outlier boundary?

\item A dataset has IQR = 20 and Q3 = 80. What is Q1?

\item A box plot shows that 50\% of the data falls between Q1 = 30 and Q3 = 50. An outlier is defined as any value more than 1.5 × IQR beyond the quartiles. What is the upper outlier boundary?

\end{enumerate}

\section{Part C: Combining Multiple Data Sources (Questions 21-30)}

\begin{enumerate}
\setcounter{enumi}{20}

\item Two stores reported their daily sales over three days:
\begin{center}
\begin{tabular}{|c|c|c|c|}
\hline
Store & Day 1 & Day 2 & Day 3 \\
\hline
Store A & \$500 & \$620 & \$580 \\
Store B & \$450 & \$510 & \$640 \\
\hline
\end{tabular}
\end{center}
Which store had higher total sales over the three days?

\item Using the data from Question 21, what was the combined total sales for both stores?

\item Two classes took a test with the following results:
\begin{center}
\begin{tabular}{|c|c|c|}
\hline
Class & Number of Students & Average Score \\
\hline
Class A & 25 & 78 \\
Class B & 30 & 82 \\
\hline
\end{tabular}
\end{center}
What is the combined average score for both classes?

\item A company has three departments with the following data:
\begin{center}
\begin{tabular}{|c|c|c|}
\hline
Department & Employees & Avg. Salary \\
\hline
Sales & 20 & \$45,000 \\
Marketing & 15 & \$55,000 \\
Operations & 35 & \$50,000 \\
\hline
\end{tabular}
\end{center}
What is the overall average salary across all departments?

\item Two gyms track membership over four months:
\begin{center}
\begin{tabular}{|c|c|c|}
\hline
Month & Gym X & Gym Y \\
\hline
Jan & 150 & 200 \\
Feb & 165 & 210 \\
Mar & 180 & 205 \\
Apr & 195 & 220 \\
\hline
\end{tabular}
\end{center}
Which gym had the greater increase from January to April?

\item Using the data from Question 25, what is the average monthly membership for Gym X over the four months?

\item Two regions reported quarterly revenue (in millions):
\begin{center}
\begin{tabular}{|c|c|c|c|c|}
\hline
Region & Q1 & Q2 & Q3 & Q4 \\
\hline
North & 12 & 15 & 18 & 22 \\
South & 10 & 14 & 16 & 20 \\
\hline
\end{tabular}
\end{center}
What is the total revenue for both regions across all quarters?

\item Using the data from Question 27, which quarter had the highest combined revenue from both regions?

\item A survey asked people in three age groups about hours of exercise per week:
\begin{center}
\begin{tabular}{|c|c|c|}
\hline
Age Group & Sample Size & Avg. Hours \\
\hline
18-30 & 120 & 5.2 \\
31-45 & 100 & 4.1 \\
46+ & 80 & 2.8 \\
\hline
\end{tabular}
\end{center}
What is the overall average hours of exercise for all respondents?

\item Using the exercise data from Question 29, how many total respondents were surveyed?

\end{enumerate}

\section{Part D: Mean, Median, and Mode (Questions 31-40)}

\begin{enumerate}
\setcounter{enumi}{30}

\item Find the mean of the dataset: 12, 15, 18, 14, 16

\item Find the median of the dataset: 24, 31, 19, 28, 25, 22, 29

\item Find the mode of the dataset: 5, 7, 5, 9, 5, 8, 7, 5

\item A dataset has the values: 8, 12, 15, 20, 25. If a value of 100 is added, how does the median change?

\item The mean of five numbers is 24. If four of the numbers are 18, 22, 26, and 28, what is the fifth number?

\item A dataset has six values with a median of 30. If the values in order are 15, 22, $x$, 35, 40, 48, what is $x$?

\item Find the mean, median, and mode for: 42, 48, 42, 45, 50, 42, 48

\item A teacher grades papers and records scores: 85, 92, 88, 90, 85, 92, 88, 95. What is the mode?

\item The mean age of a group of 8 people is 35 years. If a 9th person aged 50 joins the group, what is the new mean age?

\item A dataset has the values: 10, 15, 20, 25, 30, 35, 40. Which measure of central tendency (mean, median, or mode) best represents this data and why?

\end{enumerate}

\section{Part E: Weighted Averages and Complex Problems (Questions 41-50)}

\begin{enumerate}
\setcounter{enumi}{40}

\item A student's grade is calculated as follows: Homework (30\%), Midterm (40\%), Final (30\%). The student earned 85 on homework, 78 on the midterm, and 88 on the final. What is the weighted average grade?

\item A recipe calls for 2 cups of sugar (worth \$3 per cup) and 3 cups of flour (worth \$1 per cup). What is the average cost per cup of the final mixture?

\item A portfolio contains 100 shares of Stock A (\$50/share) and 200 shares of Stock B (\$75/share). What is the weighted average price per share?

\item An employee worked 20 hours at \$15/hour and 10 hours at \$18/hour. What is the weighted average hourly wage?

\item A company's quarterly profits are: Q1: \$50,000 (weight 1), Q2: \$60,000 (weight 2), Q3: \$55,000 (weight 1.5), Q4: \$70,000 (weight 1.5). What is the weighted average quarterly profit?

\item A store sells three types of apples at different prices:
\begin{center}
\begin{tabular}{|c|c|c|}
\hline
Type & Pounds Sold & Price per Pound \\
\hline
Red & 80 & \$1.50 \\
Green & 60 & \$1.75 \\
Golden & 40 & \$2.00 \\
\hline
\end{tabular}
\end{center}
What is the weighted average price per pound of apples sold?

\item A class with 25 students has an average test score of 82. Another class with 20 students has an average test score of 88. What is the combined average for both classes?

\item A weighted average calculation uses weights of 0.2, 0.3, 0.25, and 0.25 for values 60, 75, 80, and 90 respectively. What is the weighted average?

\item In a survey, 40\% of respondents chose Option A (satisfaction rating 8), 35\% chose Option B (satisfaction rating 6), and 25\% chose Option C (satisfaction rating 5). What is the weighted average satisfaction rating?

\item A company manufactures a product in two facilities:
\begin{center}
\begin{tabular}{|c|c|c|}
\hline
Facility & Units Produced & Cost per Unit \\
\hline
Facility A & 500 & \$8 \\
Facility B & 700 & \$6 \\
\hline
\end{tabular}
\end{center}
What is the weighted average cost per unit?

\end{enumerate}

\newpage

\section{Answer Key}
See the Solutions document for complete worked solutions to all 50 questions.

\end{document}